\chapter{Introduction}

% Background & motivation ~ why TAA and sorghum?
\hspace{1em}Both the growing concern over environmental sustainability and the depletion of fossil fuel reserves 
have driven an increasing interest in bio-based chemicals as alternatives to fossil fuel-derived compounds. 
Among the industrial sectors affected by this growing trend is the plastic manufacturing industry. 
When manufacturing a plastic product, a substance called plasticizer is incorporated to modify the 
physicochemical properties of the polymer matrix, thereby increasing flexibility and processability. 
di(2-ethylhexyl)phthalate (DEHP) is a type of phthalic acid esters (PAEs) that is known to be chemically hazardous, 
yet it is widely used as plasticizer \cite{bruni_aconitic_2022}. trans-Aconitic acid - also known as trans-aconitate or abbreviated as TAA -
is a naturally occurring secondary metabolite found in plant species such as sugar cane (\textit{Saccharum officinarum}) and sorghum 
(\textit{Sorghum bicolor}). TAA has been proposed as a bio-based precursor for plasticizer production, 
offering a sustainable alternative to synthetic PAEs \cite{bruni_aconitic_2022}.

\hspace{1em}Beyond its role as a plasticizer precursor, TAA has been identified as a versatile building block for bio-based polymer
synthesis with utility across a broad range of applications \cite{bruni_aconitic_2022}. In the biomedical field, TAA has been employed
as a constituent of scaffolds for bone tissue engineering, and has further demonstrated applicability in drug delivery systems,
where its incorporation into chitosan microparticles modulates their biocompatibility and biodegradability to improve therapeutic efficacy.

\hspace{1em}TAA can be obtained via two natural sources: microbial biosynthesis using bioengineered strains \cite{guo_efficient_2026},
or extraction from plant biomass, primarily from the stem juice of sugarcane (\textit{Saccharum officinarum}) and sweet sorghum
(\textit{Sorghum bicolor}) \cite{bruni_aconitic_2022}. Among plant-based sources, sorghum is generally regarded as the more
suitable feedstock. In sugar cane, the stem juice is the primary source of refined sugar as well as the feedstock, creating
an inherent competition between food and industrial uses. Sorghum avoids this issue: its grain is used for food, while its stem —
which has no comparable food value — can be allocated to TAA extraction independently. This separation of outputs makes sorghum
a practical and sustainable choice among the available natural sources of TAA production.

% Role of TAA in plant physiology
